\section{A Typed Temporal Semantic Representation}
\label{sec:representation}

\subsection{Schema-Guided Ontology Instantiation}

The ontology is fixed in code and specifies the allowed objects and relationships.
The generator proposes instances within this vocabulary. Formally, let
\begin{equation}
\begin{split}
 \Omega&=(\mathcal T,\mathcal R,\operatorname{dom},\operatorname{ran},\mathcal C),\\
 G_t&=(V_t,E_t,\vartheta_V,\vartheta_E,a_t).
\end{split}
\label{eq:ontology}
\end{equation}
Here $\mathcal T$ and $\mathcal R$ are the node and relation vocabularies.
$\vartheta_V:V_t\to\mathcal T$ and $\vartheta_E:E_t\to\mathcal R$
assign their types. An edge of type $r$ connects allowed types in
$\operatorname{dom}(r)$ and $\operatorname{ran}(r)$. Attributes $a_t$ record
immutable and logical IDs, dataset, subject, session, event interval, ingestion
time, version, value, unit and evidence state. Constraints $\mathcal C$ require
immutable identifiers, finite acyclic declared dependencies and the admission
checks below. Ingestion time $\tau_v$ is distinct from the type maps $\vartheta$.
At knowledge time $t$, only records with $\tau_v\leq t$ are visible.
Assessments become visible at their recorded knowledge time.
Queries select the greatest version visible at their knowledge time,
using ingestion time and ID to break ties. Historical queries retain that time. Tied assessments retain their trigger IDs. Conflicting states
remain ambiguous because the timestamp does not establish the order of triggers.

Table~\ref{tab:objects} lists the stored objects. Observation windows carry hashes,
sample indices and channel metadata, while raw arrays remain outside the database.
Features carry units, extractor metadata and source IDs. The
\texttt{with\_entities} function adds subject and channel identities.
The model proposes proposition fields, evidence references, parent links and
sentences. Claims that pass the bounded checks are inserted immutably.
The Neo4j adapter adds \texttt{SupportAssessment} nodes and
\texttt{ASSESSED\_AS} edges to the storage totals.
Inserting an edge can create a placeholder for an unresolved source Record.
None occurs in the original exports. Table~\ref{tab:relations} gives relation
meanings and creation rules. Neo4j enforces unique scope/ID pairs.
Python checks record kinds, immutable payloads, acyclicity and numerical
admission. The database itself does not enforce these semantic checks.

\begin{table}[!htbp]
\caption{Implemented object vocabulary. Checked status records an admission or
assessment operation. It does not establish physiological truth.}
\label{tab:objects}
\centering\small
\begin{tabular}{lll}
\toprule
Stored object & Meaning & Construction \\
\midrule
\shortstack[l]{Subject / Sensor} & Identity and channel & deterministic \\
Observation & Hashed sample window & deterministic \\
FeatureVersion & Value or availability watch & deterministic \\
ClaimVersion & Bounded proposition, references & model; checked \\
ExplanationVersion & Text and display membership & administrative \\
ReviewEvent & Recorded review action & administrative \\
SupportAssessment & State, time and trigger & evaluation \\
Record placeholder & Unresolved referenced ID & administrative \\
\bottomrule
\end{tabular}

\end{table}
\begin{table}[!htbp]
\caption{Relation construction and direction. R denotes a Record, P an unresolved
placeholder, F a feature, C a claim, O an observation, X an explanation and A an
assessment. The adapter constructs edges; admission checks supply claim-specific
semantics. Administrative relations do not express entailment.}
\label{tab:relations}
\centering\footnotesize\setlength{\tabcolsep}{3pt}
\begin{tabularx}{\textwidth}{l l X}
\toprule Relation & Domain $\to$ range & Meaning and creation rule\\\midrule
\texttt{DEPENDS\_ON} & R $\to$ R/P & Each source ID: deterministic input, model citation or declared parent.\\
\texttt{DERIVED\_FROM} & F $\to$ R/P & Feature source IDs: extraction provenance.\\
\texttt{SUPPORTS} & R/P $\to$ C & Reverse source edge on admission; validity at insertion only.\\
\texttt{SUPERSEDES} & R $\to$ R/P & New version to its predecessor ID.\\
\texttt{CONTAINS} & X $\to$ R/P & Explanation source IDs: display membership; also dependencies.\\
\texttt{ASSESSED\_AS} & R $\to$ A & Outcome keyed by scope, record, knowledge time and trigger.\\
\texttt{CONTRADICTS} & R $\to$ R & Trigger to adversely assessed record; time from assessment.\\
\texttt{FOR\_SUBJECT} & R $\to$ subject & Subject membership from deterministic identities.\\
\texttt{OBSERVED\_BY} & O $\to$ sensor & Channel identity from recording metadata.\\
\texttt{REVIEWED\_BY\_EVENT} & R/P $\to$ review & Reviewed item to appended review event.\\
\bottomrule
\end{tabularx}
\end{table}


Construction first normalizes the available evidence, then checks the model's
proposals and inserts the accepted records:
\begin{equation}
 G_t=\Gamma_\Omega(E_{\leq t},\widehat C_{\leq t})
 =\operatorname{Insert}_{\rm immutable}\!\left(
 \operatorname{Check}_q\!\left(\operatorname{Normalize}_\Omega(E_{\leq t}),
 \widehat C_{\leq t}\right)\right).
\label{eq:construction}
\end{equation}
$E_{\leq t}$ is the observation and revision log. $\widehat C_{\leq t}$ contains
proposals generated from the normalized visible evidence. Rejected proposals
remain in the request archive. New assessments and explanation versions extend
the graph without overwriting earlier content. The process instantiates the
specified ontology rather than discovering one.

\begin{samepage}
These rules explain much of the graph's topology. Each \texttt{source\_ids}
entry creates a \texttt{DEPENDS\_ON} edge from dependent to input.
Features also receive \texttt{DERIVED\_FROM} edges, and accepted claims receive
\texttt{SUPPORTS} edges in the opposite direction. Display membership creates
both \texttt{CONTAINS} and \texttt{DEPENDS\_ON} edges. Subject membership creates
\texttt{FOR\_SUBJECT} hubs. Our analysis separates these administrative,
provenance and reasoning roles.

\end{samepage}

\begin{figure}[!htbp]
\centering
\includegraphics[width=\textwidth]{semantic_analysis_v1/ontology_construction.pdf}
\caption{Fixed vocabulary and an actual instance. Blue denotes deterministic
recording and feature construction, green a checked model proposition, grey
identity or display administration, and amber assessment or version effects.
The concrete path comes from the exported WESAD S11 correction scope.
Its source hash, sample indices and full IDs accompany the figure.
Arrows preserve stored direction. The schema shown is a readable subset of
the complete construction table.}
\label{fig:ontology}
\end{figure}

In the saved WESAD S11 episode (Figure~\ref{fig:ontology}), the EDA window at 1275--1305 s produces a median of 2.408303 $\mu$S. The model proposes 2.408303 $\mu$S and cites that feature. The checker admits the bounded statement and insertion records both its \texttt{DEPENDS\_ON} citation and the reverse \texttt{SUPPORTS} edge. A later feature version changes the median to 2.908303 $\mu$S, retaining the original observation hash and immutable claim citation. Re-evaluation changes the claim's current support without rewriting its history.


\subsection{Edges as Explicit Semantic Commitments}

An edge's type determines what it asserts. A citation names evidence offered as support. A parent link declares a prerequisite. An extraction edge records how
a feature was produced. Membership and co-occurrence imply neither support nor
a prerequisite. In particular, an LLM-proposed \texttt{DEPENDS\_ON} edge does not
prove that its target is necessary. \texttt{SUPPORTS} records acceptance when a
claim was inserted, not validity at every later time. \texttt{CONTRADICTS}
records an adverse assessment triggered by a revision, whose knowledge time is
stored in the associated assessment. Exports retain these origins and distinguish
proposed, accepted and displayed links.

\subsection{The Query-Specific Semantic Core}

To interpret and revise an answer, a reviewer needs its claims and the inputs
they cite. We select this semantic core by a fixed projection:
\begin{equation}
 K_{q,t}=\pi_{q,\Omega}(G_t),\qquad
 V(K_{q,t})=C_{q,t}^{*}\cup F_{q,t},\quad
 E(K_{q,t})=D_{q,t}\cup P_{q,t}.
\label{eq:core}
\end{equation}
$C_{q,t}^{*}$ contains the latest displayed answer's claims and their declared
claim ancestors. $F_{q,t}$ contains the visible feature propositions they cite,
including explicit availability watches. $D_{q,t}$ retains the corresponding
stored dependency edges. When a cited feature has a newer visible version,
$P_{q,t}$ records a labelled path from the original citation through the version
chain. Each such path keeps the citation and supersession witness IDs.
It is never drawn as a directly stored Neo4j relationship.

Linked metadata retains subject and sensor identities, explanation containers,
assessments and old versions. It also retains feature-to-observation provenance,
which can be expanded in the drawing. A claim's asserted fields determine its scope, not the wrapper's query interval. Full IDs,
raw sentences, original versions and assessment times remain recoverable.
The core preserves the answer's explicit commitments through this selection rule.

We analyze three distinct representations. The complete graph contains all stored
objects. The active projection at $t$ contains the latest evidence and current
display members across queries, with their observation provenance.
The core $K_{q,t}$ answers focused questions. Which window supports
a value, and which operands support a comparison? Which version was known,
what changed after correction, and does another witness still support the claim?
We evaluate these competency questions against immutable input logs or exact
program specifications, independently of database support flags.

\subsection{Producing the Human-Review Network View}
\label{sec:reviewview}

The review view expands the semantic core for a selected answer or focus claim.
It restores original cited versions, required observation provenance and visible
successor versions instead of redirecting immutable citations.
For a stored-node budget $B$ and selected recorded times $T$, the implemented
transformation is
\begin{equation}
 H_{q,t}=\operatorname{View}_B(G_t,q),\quad
 p=\operatorname{dot}\!\left(\bigcup_{u\in T}H_{q,u}\right),\quad
 I_{q,t}=\operatorname{Render}(H_{q,t},p).
\label{eq:reviewview}
\end{equation}
$H$ is a bounded review expansion of $K$, not another learned representation.
$p$ stores node coordinates and routed edge splines. Deterministic Graphviz
layout uses sorted immutable IDs, labels and connectivity, with no random seed.
A comparison uses the union layout and reuses $p$; records not yet ingested
remain absent from the earlier drawing. Coordinates and Euclidean distances
have no semantic-similarity or neural-state interpretation.

The implementation selects database scope, answer, focus and knowledge time,
follows declared inputs, expands feature-to-observation paths on request,
and filters versions and assessments by their ingestion/knowledge times. Event
intervals describe what was measured or asserted, separately from when it was
known. Original \texttt{DEPENDS\_ON} edges remain intact; current-version
resolution retains its citation and reverse-supersession witnesses in an explicit
derived-path sidecar. Availability, support assessment and display membership
are separate fields. Of \SemAssessmentTies{} tied record/time assessment groups,
\SemAssessmentConflicts{} contain conflicting states: all outcomes and trigger
IDs remain visible as ambiguous, while membership follows the saved explanation
version. These flags and drawings never supply evaluator reference truth.

The budget counts stored records, preferring complete alternative-witness closures
when they fit. Missing records and incomplete witness groups trigger a
\emph{partial view} label and expansion control; the drawing does not claim
that omitted context is irrelevant. Subject/sensor hubs, display containers,
assessment nodes and reverse \texttt{SUPPORTS} mirrors stay in details. For
multi-input witnesses, labeled AND enclosures group feature rows: each row
maps to its stored ID, and a dashed dependency bundle enumerates its original
edges. Enclosures and bundles are rendering constructs. Between witness sets
support is OR; within each set all inputs are required.

Green/blue/grey fills distinguish claims/features/observations; outlines and
explicit words distinguish revisions, support and withdrawal. Every arrow is
labeled and preserves stored direction or identifies a grouped rendering edge.
The same typed JSON view, ID maps and Graphviz renderer produce
Figures~\ref{fig:instances} and~\ref{fig:topology} and the existing Streamlit
dashboard. Selectboxes choose a recorded case, time, focus, witness route and
node or edge; details recover full sentences, exact values/units, intervals,
versions, source IDs and assessment provenance. Highlighting changes appearance
only. Expansion exposes omitted provenance; review actions append events in an
isolated workspace. This implements inspectable external commitments, without
claiming a measured human-performance benefit or access to neural reasoning.

